% Declaração de Autoria e de Direitos
\begin{center}
Declara\c{c}\~ao de Autoria e de Direitos
\end{center}

\vspace{0.5cm}

Eu, \emph{Lucas Garcia}, autor da monografia \emph{Gest\~ao de Contratos de Energia sob Incerteza: Compara\c{c}\~ao entre Equivalente Determin\'istico (DEQ) e Programa\c{c}\~ao Din\^amica Dual Estoc\'astica (PDDE)}, subscrevo para os devidos fins, as seguintes informa\c{c}\~oes:\\
1. O autor declara que o trabalho apresentado na disciplina de Projeto de Gradua\c{c}\~ao da Escola Polit\'ecnica da UFRJ \'e de sua autoria, sendo original em forma e conte\'udo.\\
2. Excetuam-se do item 1. eventuais transcri\c{c}\~oes de texto, figuras, tabelas, conceitos e id\'eias, que identifiquem claramente a fonte original, explicitando as autoriza\c{c}\~oes obtidas dos respectivos propriet\'arios, quando necess\'arias.\\
3. O autor permite que a UFRJ, por um prazo indeterminado, efetue em qualquer m\'idia de divulga\c{c}\~ao, a publica\c{c}\~ao do trabalho acad\^emico em sua totalidade, ou em parte. Essa autoriza\c{c}\~ao n\~ao envolve \^onus de qualquer natureza \`a UFRJ, ou aos seus representantes.\\
4. O autor pode, excepcionalmente, encaminhar \`a Comiss\~ao de Projeto de Gradua\c{c}\~ao, a n\~ao divulga\c{c}\~ao do material, por um prazo m\'aximo de 01 (um) ano, improrrog\'avel, a contar da data de defesa, desde que o pedido seja justificado, e solicitado antecipadamente, por escrito, \`a Congrega\c{c}\~ao da Escola Polit\'ecnica.\\
5. O autor declara, ainda, ter a capacidade jur\'idica para a pr\'atica do presente ato, assim como ter conhecimento do teor da presente Declara\c{c}\~ao, estando ciente das san\c{c}\~oes e puni\c{c}\~oes legais, no que tange a c\'opia parcial, ou total, de obra intelectual, o que se configura como viola\c{c}\~ao do direito autoral previsto no C\'odigo Penal Brasileiro no art.184 e art.299, bem como na Lei 9.610.\\
6. O autor \'e o \'unico respons\'avel pelo conte\'udo apresentado nos trabalhos acad\^emicos publicados, n\~ao cabendo \`a UFRJ, aos seus representantes, ou ao(s) orientador(es), qualquer responsabiliza\c{c}\~ao/ indeniza\c{c}\~ao nesse sentido.\\
7. Por ser verdade, firmo a presente declara\c{c}\~ao.\\

      \vspace{0.5cm}
      \begin{flushright}
         \parbox{10cm}{
            \hrulefill
            \vspace{-.375cm}
            \centering{Lucas Garcia}
            \vspace{0.1cm}
         }
      \end{flushright}

\pagebreak

% Copyright
      \vspace{0.5cm}

UNIVERSIDADE FEDERAL DO RIO DE JANEIRO \\
Escola Polit\'ecnica - Departamento de Eletr\^onica e de Computa\c{c}\~ao \\
Centro de Tecnologia, bloco H, sala H-217, Cidade Universit\'aria \\
Rio de Janeiro - RJ \hspace{0.3cm} CEP 21949-900\\
\vspace{0.5cm}
\paragraph{}Este exemplar \'e de propriedade da Universidade Federal do Rio de Janeiro, que poder\'a inclu\'i-lo em base de dados, armazenar em computador, microfilmar ou adotar qualquer forma de arquivamento.
\paragraph{}\'E permitida a men\c{c}\~ao, reprodu\c{c}\~ao parcial ou integral e a transmiss\~ao entre bibliotecas deste trabalho, sem modifica\c{c}\~ao de seu texto, em qualquer meio que esteja ou venha a ser fixado, para pesquisa acad\^emica, coment\'arios e cita\c{c}\~oes, desde que sem finalidade comercial e que seja feita a refer\^encia bibliogr\'afica completa.
\paragraph{}Os conceitos expressos neste trabalho s\~ao de responsabilidade do(s) autor(es).

\pagebreak

% Dedicatória
\begin{center}
\textbf{DEDICAT\'ORIA}
\end{center}
      \vspace{0.5cm}

\paragraph{}% TODO: escrever dedicatória

\pagebreak

% Agradecimento
\begin{center}
\textbf{AGRADECIMENTO}
\end{center}
      \vspace{0.5cm}

\paragraph{}% TODO: escrever agradecimentos

\pagebreak

% Resumo
\begin{center}
\textbf{RESUMO}
\end{center}
      \vspace{0.5cm}

\paragraph{}Este trabalho aborda a otimiza\c{c}\~ao estoc\'astica do portf\'olio de contratos de uma comercializadora de energia el\'etrica, considerando decis\~oes de compra e venda sob incerteza e exposi\c{c}\~ao ao mercado de curto prazo. Comparam-se duas abordagens de solu\c{c}\~ao: a formula\c{c}\~ao determin\'istica equivalente (DEQ) e a Programa\c{c}\~ao Din\^amica Dual Estoc\'astica (PDDE). A avalia\c{c}\~ao, realizada por meio de experimentos computacionais com inst\^ancias sint\'eticas constru\'idas a partir de dados hist\'oricos do Sistema Interligado Nacional, indica que o DEQ apresenta limita\c{c}\~oes de escalabilidade associadas ao crescimento da \'arvore de cen\'arios e ao elevado consumo de mem\'oria. Em contrapartida, a PDDE mostra-se mais adequada para inst\^ancias de maior porte, com maior estabilidade e efici\^encia computacional. O trabalho documenta ainda a implementa\c{c}\~ao computacional completa de ambas as abordagens em linguagem Julia, com pipeline de gera\c{c}\~ao de dados em Python.

\paragraph{}
\noindent Palavras-Chave: Otimiza\c{c}\~ao Estoc\'astica, Programa\c{c}\~ao Din\^amica Dual Estoc\'astica, Mercado de Energia, Gest\~ao de Portf\'olio, Comercializa\c{c}\~ao de Energia.

\pagebreak

% Abstract
\begin{center}
\textbf{ABSTRACT}
\end{center}
      \vspace{0.5cm}

\paragraph{}This work addresses the stochastic optimization of the contract portfolio of an electricity trading company, considering purchase and sale decisions under uncertainty and exposure to the short-term market. Two solution approaches are compared: the Deterministic Equivalent (DEQ) and Stochastic Dual Dynamic Programming (SDDP). Computational experiments using synthetic instances built from historical data of the Brazilian Interconnected Power System indicate that the DEQ presents scalability limitations due to the growth of the scenario tree and high memory requirements. In contrast, SDDP is more suitable for larger-scale instances, showing greater stability and computational efficiency. The work also documents the complete computational implementation of both approaches in the Julia language, with a data generation pipeline in Python.

\paragraph{}
\noindent Key-words: Stochastic Optimization, Stochastic Dual Dynamic Programming, Energy Market, Portfolio Management, Energy Trading.

\pagebreak

% Siglas
\begin{center}
\textbf{SIGLAS}
\end{center}
      \vspace{0.5cm}

\paragraph{}ANEEL --- Ag\^encia Nacional de Energia El\'etrica
\paragraph{}CCEE --- C\^amara de Comercializa\c{c}\~ao de Energia El\'etrica
\paragraph{}CMO --- Custo Marginal de Opera\c{c}\~ao
\paragraph{}CVaR --- \textit{Conditional Value-at-Risk}
\paragraph{}DEQ --- Formula\c{c}\~ao Determin\'istica Equivalente
\paragraph{}MWm --- Megawatt-m\'edio
\paragraph{}MWh --- Megawatt-hora
\paragraph{}PAR(1) --- Modelo Autorregressivo Peri\'odico de Ordem 1
\paragraph{}PDDE --- Programa\c{c}\~ao Din\^amica Dual Estoc\'astica
\paragraph{}PLD --- Pre\c{c}o de Liquida\c{c}\~ao das Diferen\c{c}as
\paragraph{}SDDP --- \textit{Stochastic Dual Dynamic Programming}
\paragraph{}SIN --- Sistema Interligado Nacional
\paragraph{}UFRJ --- Universidade Federal do Rio de Janeiro

\pagebreak
